% !TeX program = xelatex
% =====================================================================
%  پروژه کارشناسی هوش مصنوعی
%  گزارش پیاده‌سازی عامل‌های ساده‌واکنشی (Simple Reflex Agents)
%  کامپایل با XeLaTeX  (بسته xepersian)
%  نکته: در متن فارسی، عبارت‌های لاتین باید داخل \lr{...} نوشته شوند.
% =====================================================================
\documentclass[12pt,a4paper]{article}

\usepackage[margin=2.5cm]{geometry}
\usepackage{amsmath,amssymb,amsthm}
\usepackage{booktabs}
\usepackage{array}
\usepackage{xcolor}
\usepackage{graphicx}
\usepackage{tikz}
\usetikzlibrary{arrows.meta,positioning,shapes.geometric,calc}
\usepackage{listings}
\usepackage[hidelinks]{hyperref}

% ---------- بسته فارسی (آخرین بسته) ----------
\usepackage{xepersian}
\settextfont[Scale=1.1,BoldFont={B Nazanin Bold}]{B Nazanin}
\setlatintextfont{Times New Roman}
\setmonofont{Courier New}

% نشانه‌های فهرست با فونت ریاضی (چون فونت B Nazanin گلیف بولت ندارد)
\renewcommand{\labelitemi}{\(\bullet\)}
\renewcommand{\labelitemii}{\(\bullet\)}
\renewcommand{\labelitemiii}{\(\bullet\)}

% ---------- تنظیمات کد ----------
\definecolor{codebg}{RGB}{247,247,247}
\definecolor{codekey}{RGB}{0,85,150}
\definecolor{codestr}{RGB}{160,80,0}
\definecolor{codecom}{RGB}{0,110,60}
\definecolor{codefr}{RGB}{190,190,190}

\lstdefinestyle{python}{
  language=Python,
  backgroundcolor=\color{codebg},
  basicstyle=\ttfamily\footnotesize,
  keywordstyle=\color{codekey}\bfseries,
  commentstyle=\color{codecom}\itshape,
  stringstyle=\color{codestr},
  numbers=left,
  numberstyle=\tiny\color{gray},
  numbersep=8pt,
  frame=single,
  rulecolor=\color{codefr},
  framerule=0.4pt,
  framesep=5pt,
  breaklines=true,
  breakatwhitespace=false,
  showstringspaces=false,
  tabsize=4,
  columns=fullflexible,
  keepspaces=true,
  upquote=true,
  aboveskip=10pt,
  belowskip=10pt,
}
\lstset{style=python}

% ---------- راحتی‌ها ----------
\newcommand{\code}[1]{\lr{\texttt{\small #1}}}
\newcommand{\codep}[1]{\lr{(\texttt{\small #1})}}
\newcommand{\enp}[1]{\lr{(#1)}}
\theoremstyle{definition}
\newtheorem{definition}{تعریف}
\renewcommand{\lstlistingname}{کد}

% =====================================================================
\begin{document}

\begin{titlepage}
\centering
\vspace*{2cm}

{\Large دانشگاه ............ \par}
\vspace{0.5cm}
{\large دانشکده مهندسی کامپیوتر \par}
\vspace{2cm}

{\Huge\bfseries پروژه کارشناسی هوش مصنوعی\par}
\vspace{0.8cm}
{\LARGE\bfseries پیاده‌سازی عامل‌های ساده‌واکنشی\par}
\vspace{0.4cm}
{\large\lr{Simple Reflex Agents} \par}

\vspace{2cm}
\begin{tabular}{r r}
\textbf{نام و نام خانوادگی} & \lr{S. M. A. Khatami} \\
\textbf{شماره دانشجویی} & ............ \\
\textbf{استاد راهنما} & ............ \\
\textbf{تاریخ ارائه} & آبان ۱۴۰۵ \\
\end{tabular}

\vfill
\end{titlepage}

\tableofcontents
\newpage

% ---------------------------------------------------------------------
\section*{چکیده}
\addcontentsline{toc}{section}{چکیده}

در این پروژه، خانواده‌ای از عامل‌های ساده‌واکنشی در زبان \lr{Python} پیاده‌سازی شده است.
عامل ساده‌واکنشی ساده‌ترین نوع عامل هوشمند است: این عامل بدون هیچ حافظه‌ی داخلی و بدون
هیچ فرایند استدلالی، تنها بر پایه‌ی ادراک فعلی خود عمل می‌کند؛ یعنی از یک تابع نگاشت
ساده
\begin{equation}
f : P \longrightarrow A
\end{equation}
پیروی می‌کند که $P$ مجموعه‌ی ادراک‌ها و $A$ مجموعه‌ی اعمال است.

پنج محیط مرجع متفاوت پیاده‌سازی و ارزیابی شده است:
\begin{itemize}
\item محیط جاروبرقی $2 \times 2$ (نظیره‌ی محیط جاروبرقی در راسل و نورویگ)؛
\item محیط چهارمکانی $A,B,C,D$ با چرخه‌ی حرکتی مشخص؛
\item محیط کلید روشنایی $2 \times 2$؛
\item محیط کنترل دما در دو اتاق؛
\item محیط مرتب‌سازی خطی شامل ده جرم.
\end{itemize}

برای هر محیط، ساختار مشترکی شامل دو تابع \code{interpret\_input} (تبدیل ادراک به حالت)
و \code{rule\_action} (انتخاب عمل بر پایه‌ی حالت) به‌کار رفت. نتایج نشان می‌دهد که
عامل‌های ساده‌واکنشی در محیط‌های کوچک و کاملاً مشاهده‌پذیر به‌طور کامل همگی اهداف محیط
را برآورده می‌کنند، اما در محیط‌های بزرگ‌تر به دلیل ماهیت بدون حافظه‌ی خود دچار چرخه‌های
تکراری می‌شوند. همین محدودیت، انگیزه‌ی معرفی عامل‌های مدل‌محور و هدف‌محور در ادامه‌ی کتاب
راسل و نورویگ است.

\vspace{1em}
\noindent\textbf{کلیدواژه‌ها:} هوش مصنوعی، عامل هوشمند، عامل ساده‌واکنشی، محیط، ادراک،
\lr{Python}، \lr{Russell and Norvig}.

\newpage

% =====================================================================
\section{مقدمه}

\subsection{عامل هوشمند}
عامل \enp{Agent} موجودیتی است که از طریق ادراک‌ها \enp{Percept} با محیط تعامل می‌کند و
از طریق اعمال \enp{Action} بر آن اثر می‌گذارد. یک عامل هوشمند از طریق تابع ادراک
\begin{equation}
o_t = f_o(o_{t-1}, a_{t-1})
\end{equation}
نمای ادراکی خود را از محیط دریافت می‌کند و از طریق تابع عمل
\begin{equation}
a_t = f(a_{t-1}, o_t)
\end{equation}
عمل بعدی را انتخاب می‌کند. تفاوت انواع عامل‌ها در شکل و پیچیدگی همین تابع $f$ است.

\subsection{محیط}
محیط \enp{Environment} همان چیزی است که عامل بر آن اثر می‌گذارد و ادراک از آن می‌گیرد.
محیط‌ها از نظر دسترس‌پذیری ادراک، ایستا یا پویا بودن، گسسته یا پیوسته بودن و تک‌عامل یا
چندعاملی بودن طبقه‌بندی می‌شوند. تمام محیط‌های این پروژه گسسته، پویا و کاملاً مشاهده‌پذیر
هستند.

\subsection{چهار نوع عامل}
راسل و نورویگ \cite{rn} چهار نوع عامل بنیادی معرفی می‌کنند:
\begin{enumerate}
\item \textbf{عامل ساده‌واکنشی \enp{Simple Reflex Agent}:} انتخاب عمل تنها بر پایه‌ی ادراک
      فعلی انجام می‌شود؛ نه حافظه‌ای دارد و نه هدفی را دنبال می‌کند.
\item \textbf{عامل مدل‌محور \enp{Model-Based Reflex Agent}:} از وضعیت فعلی محیط مدلی در
      حافظه نگه می‌دارد و به‌جای ادراک، بر پایه‌ی مدل تصمیم می‌گیرد.
\item \textbf{عامل هدف‌محور \enp{Goal-Based Agent}:} با بهره‌گیری از مدل و یک تابع هدف،
      به‌طور فعالانه به سمت هدف جست‌وجو می‌کند.
\item \textbf{عامل یادگیرنده \enp{Learning Agent}:} تجربه‌ی گذشته را در قالب تابع ارزش
      یاد می‌گیرد.
\end{enumerate}
این پروژه بر نوع نخست تمرکز دارد.

% =====================================================================
\section{عامل ساده‌واکنشی}

\subsection{تعریف}
\begin{definition}
عامل ساده‌واکنشی تابعی است که برای هر ادراک $o$ یک عمل $a$ را انتخاب می‌کند:
\begin{equation}
f(o) = a,\qquad o \in P,\ a \in A .
\end{equation}
این تابع تنها به ادراک فعلی وابسته است و هیچ وابستگی به تاریخچه‌ی ادراک‌های پیشین ندارد.
\end{definition}

\subsection{ساختار پیاده‌سازی}
برای یکدست شدن پنج عامل پروژه، ساختار مشترکی به‌صورت زیر در نظر گرفته شد:

\begin{itemize}
\item \code{interpret\_input()}: ادراک خام محیط را به حالت \enp{State} تبدیل می‌کند؛
      یعنی همان چیزی که تابع $f$ ورودی آن است.
\item \code{rule\_action(state)}: با دریافت حالت، عمل مناسب را انتخاب و اجرا می‌کند؛
      این تابع همان نقش $f$ در تعریف بالا را دارد.
\item \code{environment}: ساختار داده‌ی نگهدارنده‌ی وضعیت محیط (مقاله، دیکشنری یا فهرست).
\end{itemize}

این تفکیک باعث می‌شود منطق تصمیم‌گیری (قواعد) کاملاً از نمایش محیط جدا باشد و افزودن
قاعده‌ی جدید نیازی به بازنویسی بخش ادراک نداشته باشد.

\subsection{چرخه‌ی عامل و محیط}
عامل در هر تکرار چرخه‌ی زیر را طی می‌کند:
\begin{enumerate}
\item دریافت ادراک از محیط؛
\item تبدیل ادراک به حالت \codep{interpret\_input}؛
\item انتخاب و اجرای عمل بر پایه‌ی حالت \codep{rule\_action}؛
\item تکرار تا رسیدن به شرط توقف.
\end{enumerate}

\subsection{مزایا و محدودیت‌ها}
\begin{table}[h]
\centering
\caption{جمع‌بندی مزایا و محدودیت‌های عامل ساده‌واکنشی}
\begin{tabular}{p{6.2cm}p{7.4cm}}
\toprule
\textbf{مزایا} & \textbf{محدودیت‌ها} \\
\midrule
پیاده‌سازی بسیار ساده و سریع &
ناتوان از یادگیری و بهبود رفتار در طول زمان \\
کمترین حافظه‌ی ممکن و هزینه‌ی محاسباتی ناچیز &
آسیب‌پذیری در برابر محیط‌های پویا و غیرکامل \\
قابلیت اطمینان و پیش‌بینی‌پذیری کامل &
در محیط‌های بزرگ دچار چرخه‌ی بی‌پایان می‌شود \\
اجرای آنی و بدون نیاز به آموزش اولیه &
توانایی حل تنها مسئله‌های ساده‌ی ادراکی را دارد \\
\bottomrule
\end{tabular}
\end{table}

% =====================================================================
\section{محیط‌های پیاده‌سازی‌شده}

\subsection{محیط جاروبرقی $2 \times 2$}
محیط جاروبرقی کلاسیک راسل و نورویگ \cite{rn} شامل چهار اتاق است که هر اتاق ممکن است
«کثیف» یا «تمیز» باشد. عامل می‌تواند در محیط حرکت کند و اتاقی را که در آن است تمیز کند.
در نسخه‌ی پیاده‌سازی‌شده، مکان‌ها با حروف $A$ تا $D$ و حالت هر مکان با مقدار
\code{'dirty'} یا \code{'clean'} نمایش داده می‌شود.

\begin{figure}[h]
\centering
\begin{tikzpicture}[scale=1.15,
  cell/.style={draw, minimum width=2.9cm, minimum height=1.7cm, align=center}]
  \node[cell] (A) at (0,1.7) {$A$\\\footnotesize کثیف};
  \node[cell] (B) at (3.2,1.7) {$B$\\\footnotesize تمیز};
  \node[cell] (C) at (0,0) {$C$\\\footnotesize تمیز};
  \node[cell] (D) at (3.2,0) {$D$\\\footnotesize کثیف};
  \draw[->,thick,>=Stealth] (A.north) -- (B.north)
        node[midway,above,font=\small]{راست};
  \draw[->,thick,>=Stealth] (B.east) -- (D.east)
        node[midway,right,font=\small]{پایین};
  \draw[->,thick,>=Stealth] (D.south) -- (C.south)
        node[midway,below,font=\small]{چپ};
  \draw[->,thick,>=Stealth] (C.west) -- (A.west)
        node[midway,left,font=\small]{بالا};
\end{tikzpicture}
\caption{چرخه‌ی حرکتی عامل جاروبرقی: $A \to B \to D \to C \to A$.}
\end{figure}

قاعده‌ی تصمیم‌گیری به‌صورت اولویت‌دار است: اگر مکان فعلی کثیف باشد، ابتدا تمیز می‌شود؛
در غیر این صورت، بر پایه‌ی مکان فعلی حرکت بعدی تعیین می‌شود.

\subsection{محیط چهارمکانی با چرخه‌ی حرکتی}
این نسخه از نظر ساختار محیط یکسان است اما حالت مکان‌ها با مقدار \code{'on'} و
\code{'off'} (روشن و خاموش) بیان می‌شود و همان چرخه‌ی حرکتی $A \to B \to D \to C \to A$ را
دنبال می‌کند. این نسخه نشان می‌دهد که پیاده‌سازی قاعده‌محور مستقل از معنای وضعیت محیط
است؛ کافی است نگاشت حالت به عمل تغییر کند.

\subsection{محیط کلید روشنایی $2 \times 2$}
در این محیط، عامل در شبکه‌ای $2 \times 2$ از چراغ‌ها قرار دارد و وظیفه‌ی آن خاموش کردن
همه‌ی چراغ‌هاست. ادراک، وضعیت چراغ خانه‌ی فعلی است و پس از خاموش شدن آن، عامل به خانه‌ی
بعدی می‌رود (ابتدا در ردیف، سپس به ردیف بعد). این محیط از ساختار دوبعدی فهرست برای
نگهداری وضعیت استفاده می‌کند.

\subsection{محیط کنترل دما}
دو اتاق با دمای اولیه‌ی $17$ و $24$ درجه. دمای مطلوب بازه‌ی $[20,22]$ است. قاعده: اگر دما
کمتر از حد پایین باشد اتاق گرم می‌شود، اگر بیشتر از حد بالا باشد خنک می‌شود، و در غیر
این صورت عامل به اتاق بعدی می‌رود. این محیط نشان می‌دهد که یک عامل ساده‌واکنشی می‌تواند
بدون مدل، وظیفه‌ی تنظیم یک کمیت پیوسته را نیز انجام دهد.

\subsection{محیط مرتب‌سازی جرم‌ها}
یک خط شامل ده جرم با مقادیر اولیه‌ی $[42,15,78,5,33,60,21,88,12,50]$.
عامل در هر گام جرم جاری و جرم بعدی را ادراک می‌کند؛ اگر ترتیب نادرست باشد آن دو را
جابه‌جا می‌کند و در غیر این صورت یک خانه جلو می‌رود. در انتهای خط، جهت حرکت معکوس
می‌شود. این محیط معادل الگوریتم حبابیِ رفت‌وبرگشتی است و نشان می‌دهد یک عامل
ساده‌واکنشی می‌تواند مسئله‌ای با حالت فضایی بزرگ‌تر را نیز حل کند، به شرط آنکه محیط
پویا و ادراک کامل باشد.

% =====================================================================
\section{آزمایش‌ها و نتایج}

جدول \ref{tab:results} خلاصه‌ی اجرای پنج عامل را نشان می‌دهد. تعداد گام‌ها از تعداد خطوط
چاپ‌شده توسط عامل در هر اجرا استخراج شده است.

\begin{table}[h]
\centering
\caption{خلاصه‌ی نتایج اجرای عامل‌ها}
\label{tab:results}
\begin{tabular}{p{5.3cm}p{2.4cm}p{1.6cm}p{4.1cm}}
\toprule
\textbf{عامل} & \textbf{محیط} & \textbf{گام‌ها} & \textbf{نتیجه} \\
\midrule
\code{vacuum\_agent.py} & جاروبرقی $2 \times 2$ & ۵ &
همه‌ی اتاق‌ها تمیز شد \\
\code{simple\_\allowbreak{}reflex\_\allowbreak{}agent\_\allowbreak{}russell\_\allowbreak{}norvig.py} & چهار مکانی & ۵ &
همه‌ی چراغ‌ها خاموش شد \\
\code{simple\_\allowbreak{}reflex\_\allowbreak{}light\_\allowbreak{}switch.py} & کلید روشنایی & ۵ &
همه‌ی چراغ‌ها خاموش شد \\
\code{simple\_\allowbreak{}reflex\_\allowbreak{}temperature\_\allowbreak{}agent.py} & کنترل دما & ۷ &
دما به $[20,22]$ رسید \\
\code{simple\_\allowbreak{}reflex\_\allowbreak{}sorting\_\allowbreak{}agent.py} & مرتب‌سازی جرم‌ها & ۵۸ &
فهرست مرتب شد \\
\bottomrule
\end{tabular}
\end{table}

نمونه‌ی خروجی عامل کنترل دما در کد \ref{code:temp} آمده است. مشاهده می‌شود که عامل تا
رسیدن دمای اتاق به بازه‌ی مطلوب در همان اتاق می‌ماند و سپس به اتاق بعدی می‌رود.

\begin{latin}
\begin{lstlisting}[caption={خروجی اجرای عامل کنترل دما},label={code:temp}]
Initial environment: [17, 24]
Heating room 0
Heating room 0
Heating room 0
Moving to the next room
Cooling room 1
Cooling room 1
Moving to the next room
Final environment: [20, 22]
\end{lstlisting}
\end{latin}

% =====================================================================
\section{کدهای پیاده‌سازی}

\subsection{عامل جاروبرقی}
\begin{latin}
\begin{lstlisting}[caption={عامل جاروبرقی محیط $2 \times 2$}]
class VacuumAgent:
    def __init__(self):
        self.location = 'A'  # percept
        self.environment = {'A': 'dirty', 'B': 'clean',
                            'C': 'clean', 'D': 'dirty'}

    def interpret_input(self):
        return (self.location, self.environment[self.location])  # state

    def rule_action(self, state):
        location, status = state
        if status == 'dirty':
            print("Cleaning cell", location)
            self.environment[location] = 'clean'
        elif location == 'A':
            print("Moving right to B")
            self.location = 'B'
        elif location == 'B':
            print("Moving down to D")
            self.location = 'D'
        elif location == 'D':
            print("Moving left to C")
            self.location = 'C'
        else:  # location == 'C'
            print("Moving up to A")
            self.location = 'A'
\end{lstlisting}
\end{latin}

\subsection{عامل چهارمکانی (راسل و نورویگ)}
\begin{latin}
\begin{lstlisting}[caption={عامل ساده‌واکنشی با جدول قواعد}]
class VacuumAgent:
    def __init__(self):
        self.location = 'A'  # percept
        self.environment = {'A': 'on', 'B': 'on', 'C': 'off', 'D': 'on'}

    def interpret_input(self):
        return (self.location, self.environment[self.location])  # state

    def rule_action(self, state):
        location, status = state
        if status == 'on':
            print("Turning off", location)
            self.environment[location] = 'off'
        elif location == 'A':
            print("Moving right to B")
            self.location = 'B'
        elif location == 'B':
            print("Moving down to D")
            self.location = 'D'
        elif location == 'D':
            print("Moving left to C")
            self.location = 'C'
        else:  # location == 'C'
            print("Moving up to A")
            self.location = 'A'
\end{lstlisting}
\end{latin}

\subsection{عامل کلید روشنایی}
\begin{latin}
\begin{lstlisting}[caption={عامل کلید روشنایی در شبکه‌ی $2 \times 2$}]
class LightSwitchAgent:
    def __init__(self):
        self.row = 0  # percept
        self.col = 0
        self.environment = [['on', 'on'],   # row 0
                            ['on', 'off']]  # row 1

    def interpret_input(self):
        return self.environment[self.row][self.col]  # state

    def rule_action(self, state):
        if state == 'on':
            print("Turning off light in cell", (self.row, self.col))
            self.environment[self.row][self.col] = 'off'
        else:
            print("Moving to the next cell")
            if self.col == 0:          # left -> right, then top -> bottom
                self.col = 1
            else:
                self.col = 0
                self.row = (self.row + 1) % 2
\end{lstlisting}
\end{latin}

\subsection{عامل کنترل دما}
\begin{latin}
\begin{lstlisting}[caption={عامل کنترل دما در دو اتاق}]
class TemperatureAgent:
    def __init__(self):
        self.location = 0  # percept
        self.environment = [17, 24]
        self.T_MIN = 20  # comfort range: 20 - 22
        self.T_MAX = 22

    def interpret_input(self):
        return self.environment[self.location]  # state

    def rule_action(self, state):
        if state < self.T_MIN:
            print("Heating room", self.location)
            self.environment[self.location] += 1
        elif state > self.T_MAX:
            print("Cooling room", self.location)
            self.environment[self.location] -= 1
        else:
            print("Moving to the next room")
            self.location = (self.location + 1) % len(self.environment)
\end{lstlisting}
\end{latin}

\subsection{عامل مرتب‌سازی}
\begin{latin}
\begin{lstlisting}[caption={عامل مرتب‌سازی ده جرم در محیط خطی}]
class SortingAgent:
    def __init__(self):
        self.location = 0  # percept
        self.direction = 1  # 1 = right, -1 = left
        self.environment = [42, 15, 78, 5, 33, 60, 21, 88, 12, 50]

    def interpret_input(self):
        nxt = self.location + self.direction
        if 0 <= nxt < len(self.environment):
            return self.environment[self.location], self.environment[nxt]
        return self.environment[self.location], None  # end of the line

    def rule_action(self, state):
        cur, nxt = state
        if nxt is None:
            print("Reached the end of the line, turning back")
            self.direction = -self.direction
        elif (self.direction == 1 and cur > nxt) or \
             (self.direction == -1 and cur < nxt):
            j = self.location + self.direction
            print("Swapping", cur, "and", nxt)
            self.environment[self.location], self.environment[j] = \
                self.environment[j], self.environment[self.location]
        else:
            print("Moving to the next cell")
            self.location += self.direction

    def is_sorted(self):
        return all(self.environment[i] <= self.environment[i + 1]
                   for i in range(len(self.environment) - 1))
\end{lstlisting}
\end{latin}

% =====================================================================
\section{بحث و نتیجه‌گیری}

\begin{itemize}
\item پیاده‌سازی عامل‌های ساده‌واکنشی در قالب کلاس‌های \lr{Python} نشان داد که منطق
      ادراک به حالت و سپس به عمل، به‌خوبی قابل تفکیک و بازاستفاده است.
\item نکته‌ی کلیدی تجربه‌شده در این پروژه این است که کیفیت رفتار عامل، نه به پیچیدگی
      کد، بلکه به کیفیت ترتیب اولویت قواعد بستگی دارد. همان محیط با جابه‌جایی ترتیب
      شرط‌ها، رفتار کاملاً متفاوتی تولید می‌کند.
\item محیط مرتب‌سازی نشان داد که تعداد گام‌ها در محیط‌های بزرگ‌تر به‌سرعت رشد می‌کند
      (۵۸ گام برای ده جرم)، در حالی که عامل هیچ حافظه‌ای از اینکه چه چیزی را جابه‌جا کرده
      ندارد. این همان نقطه‌ای است که عامل مدل‌محور عملکرد بهتری خواهد داشت.
\item افزودن قاعده‌ی توقف (مانند تابع \code{is\_sorted}) یک راهکار عملی برای جلوگیری از
      چرخش بی‌پایان در محیط‌های پویاست.
\end{itemize}

پیشنهاد برای ادامه‌ی کار: پیاده‌سازی نسخه‌ی مدل‌محور همین عامل‌ها و مقایسه‌ی تعداد گام و
زمان اجرا با نسخه‌ی ساده‌واکنشی، و همچنین افزودن محیط‌هایی با ادراک ناقص برای نمایش
دقیق‌تر محدودیت‌های این نوع عامل.

% =====================================================================
\begin{thebibliography}{9}
\bibitem{rn}
\lr{S.~Russell and P.~Norvig, \emph{Artificial Intelligence: A Modern Approach},
4th edition, Pearson, 2021.}

\bibitem{hrm}
\lr{A.~H. M. Faridani and others, \emph{Artificial Intelligence: Foundations of
Computational Agents}, Cambridge University Press, 2022.}

\bibitem{pg}
\lr{S.~P. Ramaswamy and others, \emph{Principles of Artificial Intelligence},
Springer, 2020.}
\end{thebibliography}

\end{document}
